\documentclass{../notatki}

\title{Termodynamika i podstawy \\ fizyki statystycznej}

\begin{document}

Termodynamika to nauka, która rozważa układy makroskopowe. W nich ilość cząstek
rozważa się sumarycznie, a nie indywidualnie jak w mechanice.
$$
N \sim 6.023 \cdot 10^{23} \\
$$

\section{Różniczki}

Jeśli wiemy, że jakaś własność stanu $X$, jest funkcją stanu: $X =
X(x_1, x_2, \dots, x_n)$,
to możemy zapisać:
$$
dX = \pderiv{X}{x_1} dx_1 +
\pderiv{X}{x_2} dx_2 + \dots +
\pderiv{X}{x_n} dx_n
$$

Dla formy różniczkowej:
$$
dz(x, y) = M(x, y)dx + N(x, y)dy
$$
warunek zupełności to:
$$
\frac{\partial M}{\partial y} = \frac{\partial N}{\partial x}
$$
Jeśli jest on spełniony, to można powiedzieć, że $z(x, y)$ może być
funkcją stanu.

\begin{examplebox}
  Na przykład, założmy, funkcję stanu: $z(x, y) = x^3 + y^2$. Wtedy:
  $$
  dz = \pderiv{z}{x} dx + \pderiv{z}{y} dy = 3x^2 dx + 2y dy
  $$
  Ponieważ;
  $$
  \frac{\partial M}{\partial y} = \frac{\partial N}{\partial x} = 0
  $$
  to możemy powiedzieć, że $z(x, y)$ może być funkcją stanu.
\end{examplebox}

\section{Stany}

\begin{itemize}
  \item Izotermiczny - stała temperatura \\
    $dU = 0 \quad \partial W = - \partial Q$
  \item Izobaryczny - stałe ciśnienie \\
    $dp = 0, \quad dU = \partial Q, \quad \partial Q = dH$
  \item Izochoryczny - objętość jest stała\\
    $\partial W = 0, \quad dU = \partial Q, \quad dV = 0$
  \item Adiabatyczny - ciepło jest stałe\\
    $\partial Q = 0, \quad dU = \partial W$
\end{itemize}

Dla procesu quasi-statycznego, praca wykonana przez gaz przy zmianie
objętości:
$$
\partial W = -pdV
$$

\section{Zerowa zasada termodynamiki}

W stanie równowagi termodynamicznej nie zachodzą żadne systematyczne
zmzmiany parametrów. Nie zachodzą także żadne systematyczne przepływy
wielkości termodynamicznych.

Bycie w stanie równowagi termodynamicznej jest tranzytywne. Jeśli stworzymy
trzy zamknięte układy, obok siebie tak aby mogły przekazywać ciepło, to stan
się w nich wyrówna prędzej czy później. Dojdzie do stanu równowagi
termodynamicznej.

\begin{figure}[h]
  \centering
  \begin{tikzpicture}
    \node[draw, rectangle, minimum width=1.5cm, minimum height=1cm]
    (A) {$p_1V_1$};
    \node[draw, rectangle, minimum width=1.5cm, minimum height=1cm,
    right=0cm of A] (B) {$p_2V_2$};
    \node[draw, rectangle, minimum width=1.5cm, minimum height=1cm,
    right=0cm of B] (C) {$p_3V_3$};
    \node at (5, 1.5) {$p_1V_1 \leftrightarrow p_2V_2 \leftrightarrow p_3V_3$};

    \node at (1.5, 2.2) {diatermiczne};

    \draw[->, thick, red, dotted] (2, 2) -- ([yshift=0.5cm]B.east);
    \draw[->, thick, red, dotted] (1, 2) -- ([yshift=0.5cm]B.west);
  \end{tikzpicture}
  \caption{Układ układów w równowadze, które z czasem osiągną
  równowagę między sobą.}
\end{figure}

Rozróżniamy dwa rodzaje stanów równowagi. Pełna równowaga termodynamiczna, to
prawdziwie pełna równowaga. Nic się nie zmienia ani przepływa. Druga, pozwala
na drobne zmiany, w interesującej nas skali czasowej. Można też pominąć niektóre
parametry oczywiście.

\section{Pierwsza zasada termodynamiki}

W układzie makroskopowym, energia wewnętrzna układu zostaje zachowana. W skalach
mikroskopowych rozważamy energie kinetyczne i potencjalne wszystkich cząstek.
Fenomenologicznie rozważamy sumę wszystkich energii w układzie ($U$).
Istnieje wielkość $U$, która jest $f$ stanu.
$$
\partial W = \sum_k X_k dX_k
$$
gdzie $X_k$ są parametrami stanu.

Załóżmy, że $\Delta U = W$, czyli zmiana energii wewnętrznej jest równa pracy
wykonanej nad układem. Wtedy, po wykonaniu pracy, w której wrócimy do
stanu początkowego, mielibyśmy $W \neq 0 \nRightarrow \Delta U = 0$; oczywista
sprzeczność. Stąd, mamy że:
$$
dU = \partial W + \partial Q
$$

\begin{examplebox}
  Wiemy, że $U = U(T, V)$, a nad układem można wykonać jedynie pracę
  objętościową ($W_v$).

  \vspace{1em}
  Zatem, $W_v \Rightarrow \partial W = - pdV$. Z pierwszej zasady termodynamiki
  wiemy, że $dU = \partial W + \partial Q \Rightarrow dU = -pdV + \partial Q$.
  Możemy też rozpisać postać $dU$:
  $$
  dU = \pderiv{U}{T} dT +
  \pderiv{U}{V} dV
  $$
  To możemy przekształcić do:
  $$
  \partial Q = pdV + dU \Rightarrow \partial Q = pdV +
  \pderiv{U}{T} dT +
  \pderiv{U}{V} dV =
  \left(\frac{\partial U}{\partial T}\right)_VdT +
  \left[\pderiv{U}{T}_V + p\right] dV
  $$
  W tej formie dobrze widac, że $Q$ nie jest funkcją stanu, ale zależy od
  zmian w objętości i temperaturze.
\end{examplebox}

\section{Temperatura i ciepło}

Temperatura to parametr makroskopowy opisujący stan równowagi termodynamicznej.
Jeżeli dwa układy są w stanie równowagi termodynamicznej to ich temperatury są
równe. Jeśli połączymy dwa układy o różnych temperaturach to ich temperatury
się wyrównają.

Pojemnością cieplną nazywamy ilość ciepła $Q$ potrzebną do zmiany temperatury.
$$
\tcap_x = \pderiv{Q}{T}_x
$$
Z kolei, ciepło właściwe $C_x$ definiujemy jako:
$$
C_x = \frac{\tcap_x}{m}
$$
Dla pewnego stanu stałego ciśnienia ciśnienia, mamy:
$$
\tcap_p = \pderiv{Q}{T}_p = \pderiv{H}{T}_p
$$
Dla pewnego stanu stałej pojemności, mamy:
$$
\tcap_V = \pderiv{Q}{T}_V =
\pderiv{U}{T}_V
$$

\begin{examplebox}
  Załóżmy, izobarę, czyli: $U = U(T, V)$. Na podstawie z 1 zasady termodynamiki,
  oraz definicji ciepła właściwego, chcemy znaleźć związek między
  tymi własnościami.

  \vspace{1em}
  Z definicji ciepła właściwego, mamy:
  $$
  \tcap_V = \pderiv{Q}{T}_V =
  \pderiv{U}{T}_V \quad \tcap_P =
  \pderiv{Q}{T}_P
  $$
  Z 1 zasady termodynamiki można w tych warunkach wyprowadzić wzór na
  $\partial Q$ (\ref{examplebox:4.1}).
  $$
  \partial Q = \pderiv{U}{T}_V dT +
  \left[\pderiv{U}{V}_T + p\right]dV =
  \tcap_V dT + \left[\pderiv{U}{V}_T + p\right]dV
  $$
  $$
  \Downarrow \quad :dT_P \\
  $$
  $$
  \pderiv{Q}{T}_P = \tcap_V + \left[\pderiv{U}{V}_T + p\right]\pderiv{V}{T}_P
  $$
  $$
  \tcap_P - \tcap_V = \left[\pderiv{U}{V}_T + p\right]\pderiv{V}{T}_P
  $$

\end{examplebox}

\section{Gaz Doskonały}

$$
pV = nRT \Rightarrow V = \frac{RT}{p}
$$
gdzie $p$ to ciśnienie, $V$ to objętość, $n$ to ilość moli gazu, $R$
to stała gazowa, a $T$ to temperatura.

\begin{examplebox}
  Dla gazu $V = V(p, T)$, $dV = \pderiv{V}{p} dp +
  \pderiv{V}{T} dt$.
  $$
  dV = -\frac{RT}{p^2}dp + \frac{R}{p}dT
  $$
\end{examplebox}

Gaz doskonały spełnia:
\begin{itemize}
  \item cząstki poruszają się chaotycznie
  \item cząstki to punkty materialne $V = 0$
  \item cząsatki podlegają prawom Newtona
  \item poza momentami zderzeń, na cząstki nie działają siły
  \item zderzenia są sprężyste
\end{itemize}
zatem, powietrze w pojemniku spoczywającym o stałej objętości $V$
jest doskonałe.

\subsection{Prawa empiryczne}

Jest to zbiór praw, dla gazu doskonałego w specyficznych warunkach, z reguły,
ograniczenia warunków brzegowych.

\begin{figure}[!h]
  \centering
  \begin{tikzpicture}
    \begin{groupplot}[
        group style={
          group size=3 by 1,
          horizontal sep=1.5cm,
        },
        width=3cm, height=3cm,
        axis lines=middle,
        xlabel={$V$}, ylabel={$p$},
        ymin=0, ymax=3,
        xmin=-1, xmax=3,
        domain=0:3, samples=200,
        scale only axis,
        legend style={
          at={(-1.1,1)},
          anchor=north west,
          draw=black,
          fill=white
        },
        legend cell align=left,
        xtick=\empty, ytick=\empty
      ]

      \nextgroupplot[
        title={Izoterma, $pV = \textrm{const}$},
      ]
      \addplot[blue, thick]  {1/x};
      \addlegendentry{$p(V) = \frac{nRT}{V}$};
      \addlegendimage{red, thick}
      \addlegendentry{$p(T) = \frac{nRT}{V}$}
      \addlegendimage{green, thick}
      \addlegendentry{$V(T) = \frac{nRT}{p}$}

      \nextgroupplot[
        title={Izochora, $\frac{p}{T} = \textrm{const}$},
      ]
      \addplot[red,  thick]  {x + 1};

      \nextgroupplot[
        title={Izobara, $\frac{V}{T} = \textrm{const}$},
        xlabel={$T$}, ylabel={$V$},
      ]
      \addplot[green, thick] {x + 1};
    \end{groupplot}
  \end{tikzpicture}
\end{figure}

\subsection{Równanie adiabaty}

\begin{examplebox}
  Wyprowadźmy równanie adiabaty dla gazu doskonałego.
  $$
  pV = nRT \quad \tcap_P - \tcap_V = nR \quad dU = \partial W = -pdV
  \quad U = U(T)
  $$

  \vspace{1em}
  $$
  dU = \pderiv{U}{T}_V dT = \tcap_V dT = -pdV = -\frac{nRT}{V} dV
  \quad \tcap_V \frac{dT}{T} = - \frac{nRdV}{V}
  $$
  Całkujemy, aby wyeliminować $dT$ i $dV$.
  $$
  \tcap_V \ln T = - nRmV + C = -(\tcap_P - \tcap_V) mV + C
  $$
  Tutaj C to stała, naturalny efekt całkowania.
  $$
  \ln T ^{\tcap_V} = \ln \overline{C} V ^ {-(\tcap_P - \tcap_V)}
  $$
  $$
  T^{\tcap_V} = \overline{C} V ^ {\tcap_V - \tcap_P} \Rightarrow T =
  \overline{C}^x V ^ {\frac{\tcap_V - \tcap_P}{\tcap_V}}
  $$
\end{examplebox}

$$
T = \overline{C}^x V ^ {\frac{\tcap_V - \tcap_P}{\tcap_V}} =
\overline{C}^x V^{1 - \kappa}
$$
$$
nRT = \tilde{C} V^{1 - \kappa}
$$
$$
pV^{\kappa} = \tilde{C}
$$
Stałe $\overline{C}$, $\overline{C}^x$, $\tilde{C}$ to po prostu różne stałe,
wynikające z logarytmów. Zazwyczaj, te stałe są stałe dla
\textbf{danej adiabaty}. To oznacza, że można przyrównać te stałe do siebie
wraz z zmianą jakiegoś parametru (np. objętości).

\begin{examplebox}
  Zakładając środowisko adiabatyczne, trzeba wyprowadzić pracę przy zmianie
  objętości.
  $$
  \stackrel{p_1}{V_1} \to \stackrel{p_2}{V_2} \quad W = ?
  $$
  $$
  \partial W = -pdV \quad pV^\kappa = \tilde{C} \Rightarrow
  p_1V_1^\kappa = p_2V_2^\kappa \quad pV = nRT
  $$

  \vspace{1em}
  $$
  p_1V_1^\kappa = p_2V_2^\kappa \quad W = \int \partial W = -
  \int_{V_1}^{V_2} p \cdot dV = - \int_{V_1}^{V_2}
  \frac{p_1V_1^\kappa}{V^\kappa} \cdot dV = - p_1V_1^\kappa \int_{V_1}^{V_2}
  \frac{1}{V^\kappa} \cdot dV
  $$
  $$
  W = -p_1V_1^\kappa \left[\frac{V^{-\kappa + 1}}{1 -
  \kappa}\right]^{V_2}_{V_1} = \frac{p_1V_1}{\kappa - 1}
  \left(\left(\frac{V_2}{V_1}\right)^{1 - \kappa} - 1 \right)
  $$
  To jest jakaś postać, ale da się jeszcze uprościć ją.
  $$
  \left.
  \begin{aligned}
    V = \frac{nRT}{p} &\Rightarrow& \frac{V_2}{V_1} = \frac{T_2}{T_1}
    \frac{p_1}{p_2} \\
    p_1V_1^\kappa = p_2V_2^\kappa &\Rightarrow& \frac{p_2}{p_1} =
    \left(\frac{V_1}{V_2}\right)^\kappa
  \end{aligned}
\right\} \left(\frac{V_2}{V_1}\right)^{1 - \kappa} =
\left(\frac{V_2}{V_1}\right)^{1} \left(\frac{V_1}{V_2}\right)^\kappa =
\frac{T_2}{T_1} \frac{p_1}{p_2} \cdot  \frac{p_2}{p_1} = \frac{T_2}{T_1}
$$
$$
W = \frac{p_1V_1}{\kappa - 1}\left(\frac{T_2}{T_1} - 1\right) =
\frac{p_1V_1}{(\kappa - 1)T_1} (T_2 - T_1) =
\frac{nR}{\kappa - 1}(T_2 - T_1)
$$
\end{examplebox}

\subsection{Inne równania}

\begin{itemize}
\item $W = p(V_1 - V_2) = nR(T_1 - T_2)$ - izobara
\item $W = nRT \ln \frac{V_1}{V_2} = nRT \ln \frac{p_2}{p_1}$ - izoterma
\item $W = 0$ - izochora
\end{itemize}

\subsection{Doświadczenie Joule'a}

\begin{figure}[h]
\centering
\begin{tikzpicture}
  \node[draw, minimum width=1.5cm, minimum height=3cm] (A) at (0,0)
  {\textit{gaz}};
  \node[draw, minimum width=1cm, minimum height=0.3cm, right=of A,
  xshift=-1cm] (B) {};
  \node[draw, minimum width=1.5cm, minimum height=3cm, right=of B,
  xshift=-1cm] (C) {\textit{próżnia}};

  \node[align=center, yshift=2cm] at (B.north) (Btext) {przegroda};

  \node[draw, thick, rounded corners, fit=(A) (B) (C) (Btext), inner sep=8mm]
  (bigbox) {};

  \node[align=center, yshift=-0.2cm] at (bigbox.south) (bordertext)
  {osłona izolująca};

  \draw[->] (Btext.south) -- (B.north);
\end{tikzpicture}
\caption{Diagram eksperymentu Joule'a}
\end{figure}

Po usunięciu przegrody, temperatura gazu się nie zmienia. Stąd:
$$
U \neq U(V)
$$
$U$ jest funkcją temperatury.
$$
dU = n\pderiv{U}{T}_V dT = n\tcap_V dT \quad U(T) = n\tcap_V T
$$

\subsection{Rozkład Maxwella}

Rozkład pradopodobieństwa prędkości cząstek gazu doskonałego $v$ w
danej temperaturze to:
$$
\rho(v) = 4 \pi \left(\frac{m_0}{2\pi k_BT}\right)^{\frac{3}{2}} v^2
e^{-\frac{m_0v^2}{2k_BT}}
$$

\begin{examplebox}
$$
\langle f(v) \rangle = \int_0^\infty f(v) \rho(v) dv
$$

\vspace{1em}
$$
\langle v^2 \rangle = \int_0^\infty v^2 \rho(v) dv = 4 \pi
\left(\frac{m_0}{2\pi k_BT}\right)^{\frac{3}{2}} \int_0^\infty v^4
e^{-\frac{m_0v^2}{2k_BT}} dv = 4 \pi \left(\frac{m_0}{2\pi
k_BT}\right)^{\frac{3}{2}} \cdot \frac{3}{8}
\sqrt{\frac{\pi}{\frac{m_0}{2k_BT}^5}} =
$$
$$
= \frac{3}{2} \pi^{\frac{3}{2}} \left(\frac{m_0}{2\pi
k_BT}\right)^{\frac{3}{2}}
\left(\frac{m_0}{2k_BT}\right)^{-\frac{5}{2}} = \frac{3}{2}
\left(\frac{m_0}{2k_BT}\right)^{-\frac{2}{2}} = 3 \frac{k_BT}{m_0}
$$
\end{examplebox}

\subsection{Zasada ekwipartycji energii}

Na każdy stopień swobody, średnia energia kinetyczna gazu wynosi
$\frac{1}{2} k_BT$

\section{Cykle termodynamiczne}

\begin{figure}[!tbh]
\centering
\begin{tikzpicture}
  \node (a) at (1,1) {A};
  \node (b) at (1,3) {B};
  \node (c) at (3,1) {C};

  \draw[->] (-0.2,0) -- (5,0) node[right] {$V$};
  \draw (1,0.08) -- (1,-0.08) node[below] {\small $V_1$};
  \draw (3,0.08) -- (3,-0.08) node[below] {\small $V_2$};
  \draw[->] (-0.2,0) -- (-0.2,4) node[above] {$p$};
  \draw (-0.25,1) -- (-0.15,1) node[left] {\small $p_1$};
  \draw (-0.25,3) -- (-0.15,3) node[left] {\small $p_2$};

  \draw[->] (a) -- node[midway,left] {\tiny izochora} (b);
  \draw[->] (a) -- node[midway,below] {\tiny izobara} (c);
  \draw[->] (b) -- (c);

\end{tikzpicture}
\caption{Przykładowy cykl termodynamiczny}
\end{figure}

\begin{examplebox}
\begin{itemize}
  \item $A \rightarrow B$:\\
    $\partial W_{AB} = -pdV = 0$\\
    $W_{AB} = 0$
  \item $B \rightarrow C$:\\
    $\partial W_{BC} = -pdV$
    $$
    W_{BC} = - \int pdV = - \int_{V_1}^{V_2} p(v)dV
    $$
    Jeśli założymy, że $p(v)$ jest liniowa w $[V_1, V_2]$:
    $$
    W_{BC} = - \int_{V_1}^{V_2} \left(\frac{p_2 - p_1}{V_1 - V_2}(V_2 - V)
    + p_1\right)dV = -\left(p_1(V_2 - V_1) + \frac{1}{2}(p_2 -
    p_1)(V_2 - V_1)\right)
    $$
  \item $A \rightarrow C$:\\
    $\partial W_{AC} = -p_1dV$
    $$
    W_{AC} = -p_1\int_{V_1}^{V_2} dV = -p_1(V_2 - V_1)
    $$
\end{itemize}
\end{examplebox}

\subsection{Cykl Carnota}

Jest to cykl z czterema przejściami, symetrycznie, izoterma($T_1$) -
adiabata - izoterma($T_2$) - adiabata. $T_2 > T_1$. W ramach cyklu wydzielane
i pobierane jest ciepło.

\begin{figure}[!tbh]
\centering
\begin{tikzpicture}
  \node (a) at (1,5) {A};
  \node (b) at (4,4) {B};
  \node (c) at (5,1) {C};
  \node (d) at (2, 2) {D};

  \node (center) at (3,3) {$\overline{W}$};

  \draw[->] (-0.2,0) -- (6,0) node[right] {$V$};
  \draw (1,0.08) -- (1,-0.08) node[below] {\small $V_1$};
  \draw (5,0.08) -- (5,-0.08) node[below] {\small $V_2$};
  \draw[->] (-0.2,0) -- (-0.2,6) node[above] {$p$};
  \draw (-0.25,1) -- (-0.15,1) node[left] {\small $p_1$};
  \draw (-0.25,5) -- (-0.15,5) node[left] {\small $p_2$};

  \draw[->] (a) to[out=-30, in=180]
  node[midway,above,sloped] {$\stackrel{Q_{AB}}{\Downarrow}$}
  node[midway,below,sloped] {\tiny izoterma($T_2$)}(b);
  \draw[->] (b) to[out=-80, in=110] node[midway,above, sloped] {\tiny
  adiabata} (c);
  \draw[->] (c) to[out=180, in=-30]
  node[midway,below, sloped] {$\stackrel{\Downarrow}{Q_{CD}}$}
  node[midway,above,sloped] {\tiny izoterma($T_1$)} (d);
  \draw[->] (d) to[out=110, in=-80] (a);

\end{tikzpicture}
\caption{Cykl Carnota}
\end{figure}

Z definicji:
$$
\eta = \frac{\overline{W}}{Q_{AB}} = 1 - \frac{T_1}{T_2}
$$
Z pierwszej zasady termodynamiki:
$$
\overline{W} = Q_{AB} + Q_{CD}
$$

\subsection{Cykl Otta}

\begin{figure}[!tbh]
\centering
\begin{tikzpicture}
  \node (a) at (2,4) {A};
  \node (b) at (5,3) {B};
  \node (c) at (5,1) {C};
  \node (d) at (2, 2) {D};

  \draw[->] (-0.2,0) -- (6,0) node[right] {$V$};
  \draw (2,0.08) -- (2,-0.08) node[below] {\small $V_1$};
  \draw (5,0.08) -- (5,-0.08) node[below] {\small $V_2$};
  \draw[->] (-0.2,0) -- (-0.2,5) node[above] {$p$};
  \draw (-0.25,2) -- (-0.15,2) node[left] {\small $p_1$};
  \draw (-0.25,4) -- (-0.15,4) node[left] {\small $p_2$};

  \draw[->] (a) to[out=-30, in=180] node[midway,above, sloped] {\tiny
  adiabata} (b);
  \draw[->] (b) -- node[midway,right] {$\Rightarrow Q_{BC}$} (c);
  \draw[->] (c) to[out=180, in=-30] node[midway,below, sloped] {\tiny
  adiabata} (d);
  \draw[->] (d) -- node[midway,left] {$Q_{DA} \Rightarrow $} (a);

\end{tikzpicture}
\caption{Cykl Otta}
\end{figure}

$$
Q_{DA} = n \tcap_V (T_A - T_D) > 0 \quad Q_{BC} = n \tcap_V (T_C - T_B) > 0
$$
$$
\eta = 1 - \frac{T_B - T_C}{T_A - T_B} = 1 - r^{1 - \kappa}
$$
$$
r = \frac{V_C}{V_D} > 1
$$

\subsection{Cykl Diesel'a}

\begin{figure}[!tbh]
\centering
\begin{tikzpicture}
  \node (a) at (1,5) {A};
  \node (b) at (3,5) {B};
  \node (c) at (6,3) {C};
  \node (d) at (6, 1) {D};

  \draw[->] (-0.2,0) -- (7,0) node[right] {$V$};
  \draw (1,0.08) -- (1,-0.08) node[below] {\small $V_1$};
  \draw (6,0.08) -- (6,-0.08) node[below] {\small $V_2$};
  \draw[->] (-0.2,0) -- (-0.2,6) node[above] {$p$};
  \draw (-0.25,1) -- (-0.15,1) node[left] {\small $p_1$};
  \draw (-0.25,5) -- (-0.15,5) node[left] {\small $p_2$};

  \draw[->] (a) -- node[midway,below] {\tiny izobara}
  node[midway,above] {$\stackrel{Q_{AB}}{\Downarrow}$} (b);
  \draw[->] (b) to[out=-70, in=180] node[midway,above, sloped] {\tiny
  adiabata} (c);
  \draw[->] (c) -- node[midway,left] {\tiny izochora}
  node[midway,right] {$\Rightarrow Q_{CD}$} (d);
  \draw[->] (d) to[out=180, in=-70] node[midway,below, sloped] {\tiny
  adiabata} (a);

\end{tikzpicture}
\caption{Cykl Diesel'a}
\end{figure}

$$
r = \frac{V_D}{V_A} \quad k = \frac{V_B}{V_A}
$$
$$
Q_{CD} = n\tcap_V(T_D - T_C) < 0 \quad Q_{AB} = n\tcap_p (T_B - T_A) > 0
$$
$$
\eta = 1 - \frac{\tcap_V(T_C - T_D)}{\tcap_p (T_B - T_A)} = 1 -
\frac{(T_C - T_D)}{\kappa (T_B - T_A)}
$$

\section{Entropia}

Entropia to funkcja stanu, która się zmienia, gdy ciepło jest
wymieniane między układem i jego otoczeniem:
$$
dS = \frac{\partial Q}{T}
$$
$$
S = k_B \ln T
$$

\begin{examplebox}
Obliczmy entropię gazu doskonałego.
$$
S = S(V, T) \quad dU = \partial Q + \partial W = \partial Q - pdV =
n\tcap_V dT \quad pV = nRT
$$

\vspace{1em}
$$
\partial Q = pdV + n\tcap_V dT \Rightarrow dS = \frac{p}{T} dV +
\frac{n\tcap_V}{T} dT = \frac{nR}{V} dV + \frac{n\tcap_V}{T} dT
$$
Mamy teraz postać różniczkową. Aby uzyskać postać funkcji, musimy zcałkować:
$$
S_2 - S_1 = \int_1^2 dS = nR\ln\frac{V_2}{V_1} + n\tcap_V \ln\frac{T_2}{T_1}
$$
$$
S(T, V) = \int dS = nR\ln(V) + n\tcap_V\ln(T) + C(n)
$$
Jeśli wykorzystamy wzór na $\tcap_V$ dla gazu doskonałego, możemy ten
wzór jeszcze uprościć
$$
\tcap_V = \frac{\sigma}{2}R \Rightarrow S(V, T) =
nR\ln\left(VT^{\frac{\sigma}{2}}\right) + C
$$
\end{examplebox}

\subsection{Własności}

\begin{itemize}
\item Addytywność: $A = B \cup C \rightarrow S_A = S_B + S_C$
\item $dE = TdS$
\item $\partial Q = TdS$
\end{itemize}

\subsection{Entropia gazu doskonałego}

$$
S(V, T) = C_0(N) + Nk_B \ln (BT^{\frac{3}{2}})
$$

Dla 1 atmosfery mamy:
$$
n = \frac{N}{V} \quad \lambda = \sqrt{\frac{2\pi h^2}{mk_BT}} \quad S
= Nk \left[\frac{5}{2} - \ln \left(n \lambda^3\right)\right]
$$

\section{Entalpia}

Entalpia to suma wewnętrznej energii systemu termodynamicznego.
$$
H(S, p) = U + pV
$$
Jest to funkcja stanu, zależna od entropii i ciśnienia.

\section{Relacje Maxwella}

Jest to zbiór relacji, które można wyprowadzić z definicji
potencjałów termodynamicznych, przy pewnych założeniach. W ramach
relacji, zakładamy proces odwracalny na gazie doskonałym, co tłumaczy się do:
$$
\left.
\begin{aligned}
\partial Q &=& TdS\\
\partial W &=& -pdV
\end{aligned}
\right\} \Rightarrow dU = \partial Q + \partial W = TdS - pdV
$$

\begin{examplebox}
Na podstawie definicji energii wewnętrznej układu $U$, możemy wyprowadzić
pierwszą relację Maxwella.
$$
U = U(S, V)
$$

\vspace{1em}
Ponieważ:
$$
dU = \pderiv{U}{S}_V dS + \pderiv{U}{V}_S + dV = TdS - pdV
$$
to możemy zauważyć że:
$$
\pderiv{U}{S}_V = T \quad \pderiv{U}{V}_S = -p
$$
Jeśli potem obliczymy pochodne mieszane, i założymy, że muszą być
równe bo $U$ jest "ładną funkcją":
$$
\left.
\begin{aligned}
\frac{\partial}{\partial V} \pderiv{U}{S}_V &=& \pderiv{T}{V}_S\\
\frac{\partial}{\partial S} \pderiv{U}{V}_V &=& -\pderiv{p}{S}_V
\end{aligned}
\right\} \pderiv{T}{V}_S = -\pderiv{p}{S}_V
$$
\end{examplebox}

\begin{examplebox}
Na podstawie defincji entalpii, możemy wyprowadzić drugą relację Maxwella.
$$
H = U + pV = H(S, p)
$$

\vspace{1em}
Rozwijając $dH$ na podstawie definicji $H$:
$$
dH = dU + d(pV) = \partial Q + \partial W + d(pV)= TdS - pdV + Vdp +
pdV = TdS + Vdp
$$
równocześnie, z definicji $H = H(s, p)$;
$$
dH = \pderiv{H}{S}_pdS + \pderiv{H}{p}_Sdp
$$
Teraz, podobnie jak w poprzednim przykładzie \ref{examplebox:10.1},
możemy zauważyć, że:
$$
T = \pderiv{H}{S}_p \quad V = \pderiv{H}{p}_S
$$
i co za tym idzie, zakładając, że $H$ to "ładna funkcja":
$$
\pderiv{T}{p}_S = \pderiv{V}{S}_P
$$
\end{examplebox}

\begin{examplebox}
Na podstawie energii swobodnej (Helmholtza) $F$, możemy wyprowadzić kolejną
relację Maxwella.
$$
F = U - TS = F(T, V)
$$

\vspace{1em}
Rozwijając z definicji $dF$:
$$
dF = dU - SdT - TdS = TdS - pdV -SdT - TdS = -pdV - SdT
$$
Na podstawie definicji, podobnie jak w poprzednich przykładach:
$$
\pderiv{F}{T}_V = -S \quad \pderiv{F}{V}_T = -p
$$
I teraz, po podstawieniu i ponownym założeniu "ładnego" $F$:
$$
\pderiv{S}{V}_T = \pderiv{p}{T}_V
$$
\end{examplebox}

\begin{examplebox}
Na podstawie entalpii swobodnej (Gibbsa) $G$, możemy wyprowadzić
kolejną relację
Maxwella
$$
G = H - TS = G(T, p)
$$

\vspace{1em}
Rozwijając z definicji $G$:
$$
dG = dH - SdT - TdS = dU + d(pV) - SdT - TdS = \partial Q + \partial W + Vdp +
pdV - SdT - TdS =
$$
$$
= Vdp - SdT
$$
Dostajemy następujące relacje:
$$
\pderiv{G}{T}_p = -S \quad \pderiv{G}{p}_T = V
$$
które po założeniu "ładnego" $G$:
$$
-\pderiv{S}{p}_T = \pderiv{V}{T}_p
$$
\end{examplebox}

\textbf{Relacje są prawdziwe, jeśli spełnione są założenia.}

\section{Zespół kanoniczny}
% ładna nazwa nadana przez GPT, cholera wie czy ma sens

\begin{examplebox}
Dowód
$$
U = \frac{\partial}{\partial \beta} (\beta F)_V \quad \beta = \frac{1}{T}
$$

\vspace{1em}
Najpierw, rozbijemy prawą część równania:
$$
U = \frac{\partial}{\partial \beta} (\beta F)_V = F +
\pderiv{F}{\beta}_V  \beta = F + \textcolor{red}{\pderiv{F}{T}_V}
\textcolor{blue}{\pderiv{T}{\beta}_V} \beta
$$
Jeśli rozwiniemy definicję $F$, oraz jej różniczki, to dostaniemy
uproszczenie:
$$
F=U-TS \Rightarrow dF = dU - TdS - SdT = \pderiv{F}{T}_V dT +
\pderiv{F}{V}_T dV \Rightarrow \textcolor{red}{\pderiv{F}{T}_V} = -S
$$
Następnie z własności pochodnej możemy uprościć drugi czynnik:
$$
\textcolor{blue}{\pderiv{T}{\beta}_V} = \pderiv{\beta^{-1}}{\beta}_v
= - \frac{1}{\beta ^2} = -T^2
$$
Co finalnie daje nam:
$$
U = F + (-S)(-T^2)\beta = F + ST = U - TS + TS = U \quad \square
$$
\end{examplebox}

\begin{examplebox}
Dowód
$$
H = \frac{\partial}{\partial \beta} (BG)_p
$$

\vspace{1em}
Rozwijając równanie, uzyskamy znajomą nam już postać:
$$
H = \frac{\partial}{\partial \beta} (BG)_p = G + \beta \pderiv{G}{\beta}_p =
G +
\textcolor{red}{\pderiv{G}{T}_p}\textcolor{blue}{\pderiv{T}{\beta}_p} \beta
$$
Jeśli rozwiniemy definicję $G$ i jej różniczki, to dostaniemy uproszczenie:
$$
G = H - TS \Rightarrow dG = dH - TdS - SdT = -SdT + Vdp =
\pderiv{G}{T}_p dT + \pderiv{G}{p}_T dp \Rightarrow
\textcolor{red}{\pderiv{G}{T}_p} = -S
$$
Z poprzedniego przykładu już znamy drugie uproszczenie, stąd:
$$
H = G + (-S)(-T^2)\beta = G + TS = H - TS + TS = H \quad \square
$$
\end{examplebox}

\begin{examplebox}
Dowód
$$
\pderiv{\tcap_V}{V}_T = T \left(\frac{\partial^2 p}{\partial T^2}\right)_V
$$

\vspace{1em}
Jeśli rozwieniemy definicję $\tcap_V$, to dostaniemy wyrażenie z $S$:
$$
\tcap_V = \pderiv{Q}{T}_V = \left(\frac{TdS}{dT}\right)_V = T\pderiv{S}{T}_V
$$
Z definicji $S$, możemy wyrazić $\tcap_V$ przy pomocy $F$:
$$
S = -\pderiv{F}{T}_V \Rightarrow dS = - \frac{\partial}{\partial T}
\pderiv{F}{T}_V = -\left(\frac{\partial^2 F}{\partial T^2}\right)
\Rightarrow \tcap_V = -T\left(\frac{\partial^2 F}{\partial T^2}\right)
$$
Z definicji $p$, możemy prawą stronę równania zapisać też przy pomocy $F$:
$$
p = - \pderiv{F}{V}_T \Rightarrow \left(\frac{\partial^2 p}{\partial
T^2}\right)_V = -\frac{\partial^3 F}{\partial T^2 \partial V}
$$
Na koniec, z definicji różniczki, możemy przyrównać obydwie strony:
$$
\pderiv{\tcap_V}{V}_T = -T\frac{\partial}{\partial V}
\left(\frac{\partial^2 F}{\partial T^2}\right) = T\left(-\frac{\partial^3
F}{\partial T^2 \partial V}\right) = T\left(\frac{\partial^2 p}{\partial
T^2}\right)_V \quad \square
$$
\end{examplebox}

\begin{examplebox}
Dowód
$$
\pderiv{\tcap_p}{p}_T = -T\left(\frac{\partial^2 V}{\partial T^2}\right)_p
$$

\vspace{1em}
Jeśli rozwiniemy definicję $G$ i jej różniczki, to będziemy mogli
wyrazić $V$ przez $G$:
$$
G = H - TS \Rightarrow dG = dH - TdS - SdT = -SdT + Vdp =
\pderiv{G}{T}_p dT + \pderiv{G}{p}_T dp \Rightarrow \pderiv{G}{p}_T = V
$$
Teraz możemy wyrazić prawą stronę równania przez $G$:
$$
-T\left(\frac{\partial^2 V}{\partial T^2}\right)_p = -T
\frac{\partial^2}{\partial T^2} V = -T
\frac{\partial^2}{\partial T^2} \pderiv{G}{p}
$$
Na koniec wystarczy tylko rozwinąć $\tcap_p \to \partial Q \to dS \to
G$ i przyrównać obydwie strony równania:
$$
\pderiv{\tcap_p}{p}_T = \frac{\partial}{\partial p} \pderiv{Q}{T}_p =
\frac{\partial}{\partial p} \left(\frac{TdS}{dT}\right)_p = -T
\frac{\partial}{\partial p}
\left(\frac{\partial^2 G}{\partial
T^2}\right)_p = -T \frac{\partial^3 G}{\partial T^2 \partial p} \quad \square
$$
\end{examplebox}

\section{Potencjał chemiczny}

Potencjał chemiczny zapisujemy jako $\mu$.
$$
\begin{aligned}
dU = \pderiv{U}{S}_{V, n} dS + \pderiv{U}{V}_{S, n} dV +
\pderiv{U}{n}_{S, V} dn &\rightarrow& \mu = \pderiv{U}{n}_{S, V} \\
dF = -SdT - pdV + \mu dn &\rightarrow& \mu = \pderiv{F}{n}_{T, V} \\
dH = TdS + Vdp + \mu dn &\rightarrow& \mu = \pderiv{H}{n}_{s, p} \\
dG & \rightarrow & \mu = \pderiv{G}{n}_{T,P}
\end{aligned}
$$

\subsection{Wielki potencjał termodynamiczny}

$$
\Omega = F - \mu n
$$
$$
d\Omega = -SdT - pdV - nd\mu
$$

\end{document}
